\chapter{The Type C arguments}
\label{chap:current-type-c}

Let $F$ be a finite field of order $q=p^f$, let $n\geq2$, and set
$G=\Sp_{2n}(F)$ and $S=\PSp_{2n}(F)$. The pair $(n,q)=(2,2)$ is excluded.
The argument separates the defining characteristic from the nondefining
primes, then treats the parity of $q$ and the small ranks.

For $\ell=2$ with $q$ odd, principal weights are parametrised by their subgroup
factors and irreducible characters of the corresponding normaliser quotients.
Over an even field, the correspondence between
unipotent characters and generic weights treats the unipotent blocks needed for Jordan reduction. For odd nondefining primes over odd fields, an integral basic set
and Conlon's theorem give the bijection between Brauer characters and ordinary characters needed for the stabiliser argument. Jordan reduction uses a suitable automorphism group for
each semisimple parameter.

Throughout, the hypotheses identify the finite groups, their blocks and
characters, and the common roots used for ordinary and Brauer character
values. The distinctions between these interpretations and the proved
deductions are set out in Chapter~\ref{chap:current-scope}.

\section{The principal block in odd characteristic}

At rank two, manuscript Proposition 3.3 uses
\cite[Theorem~1.1]{LiLi2019}. For odd $q_0$, this gives the inductive
condition for $\PSp_4(q_0)$ at $2$. Manuscript Corollary 2.5 gives an
iBAW bijection on its principal block. Manuscript Lemma 2.6 lifts this
bijection to the principal block of $\Sp_4(q_0)$ through its central
subgroup of order two.
This supplies the rank two principal block required by Jordan reduction.
The construction below applies in rank at least three.

We use the principal weight criterion in \cite[Lemma~2.3]{FengYuZhang2024}. For
a $2$-subgroup $R$ and an irreducible character $\varphi$ of $N_G(R)$, the
criterion concerns the same character in three conditions: $Z(R)$ is a Sylow
$2$-subgroup of $C_G(R)$, the subgroup $RC_G(R)$ is contained in
$\ker\varphi$, and the character on $N_G(R)/RC_G(R)$ has defect zero.

Radicality follows from the local character conditions. A quotient with an
irreducible ordinary character of defect zero has no nontrivial normal
$2$-subgroup.
For the constructed subgroup,
$C_G(R)\leq R$, so $RC_G(R)=R$ and $C_G(R)=Z(R)$. The same quotient
character then satisfies the principal criterion. This proves that the
constructed pair is a principal weight.

Proposition 3.48 of \cite{FengYuZhang2024} lists the possible basic
factors. The manuscript excludes case (4) by a reflection argument
and uses Remark 3.50 for cases (1), (2) and (5).
The formalisation assumes these three exclusions and derives the
exclusion of case (4) from the reflection argument.
For each irreducible character of the relevant normaliser quotient,
its multiplicity is the size of a staircase partition.
Each such multiplicity is a triangular number, but their sum need
not be a triangular number.
Manuscript Remark 3.4 illustrates this distinction with
$Q_8\times Q_8$ in $\Sp_4(3)$. The Lean results
\lean{twoColouredHeights_size} and \lean{triangular_ne_two} in
\module{ModularRep.PaperProofs.OddTwoWreathCoreCharacterSource}
prove the arithmetic example with heights $(1,1,0)$. They do not
identify its symplectic subgroup or normaliser quotient.

On each symplectic summand, the centraliser calculation is the applicable
part of An's argument \cite[p.~198]{An1993}. The independent scalar
transformations which act as $-I$ on one summand and as $I$ on the others
belong to the constructed radical subgroup. Commutation with these signs
forces each matrix entry between distinct summands to equal its negative. It
vanishes in odd characteristic. The symplectic form equation then restricts
to each diagonal block, so an element centralising the whole subgroup
preserves every symplectic summand. Combining this preservation of the
summands with the local centraliser calculation gives the centraliser
containment for the direct product.

The remaining radical subgroup factors attached to the polynomial $x-1$
and their irreducible quotient characters are identified with the
construction used in Feng--Malle's Proposition 5.4
\cite{FengMalle2022}. A parameter consists of a subgroup decomposition
and an independent staircase parameter for each quotient character. A
separate hypothesis classifies these subgroup decompositions up to
conjugacy. The character parametrisation concerns the quotients of those
subgroups. The resulting maps in the two directions are proved to be inverse.

For the inverse map, assume a factor decomposition of the given principal
weight, the subgroup formula and the normalisation of its basic factors.
Together with the exclusions, these determine the surviving factors, their
symplectic summands and the subgroup identity. No description of the
normaliser or enumeration of quotient characters is assumed at this step.
The independent signs and the local centraliser calculation above prove
centraliser containment for the constructed subgroup. This containment is
then used to obtain the normaliser description.

The numerical parameters satisfy the weighted rank condition. Their colours
and diagonal action are those in Feng--Malle's local formulae and parameter
action \cite[Lemma~5.1 and Proposition~5.9]{FengMalle2022}. The comparison
maps preserve the specified subgroup, character and conjugator. Thus the
field and diagonal actions on the numerical parameters give the actions on
principal weight classes.

For principal Brauer characters in rank at least three, manuscript
Lemma 2.12 gives the ordinary character selection in
$H=\mathbf H^{F_0}=\CSp_{2m}(q_0)$, where $q_0$ is odd.
Choose the geometric semisimple element $s_0$ and Weyl group family
$\mathcal F$ in \cite[10.1--10.3]{Taylor2014Maximal}, with $s_0$
in a maximal torus of $\mathbf H^*$ split over $\mathbb F_{q_0}$.
Its image in the adjoint group is quasi-isolated
\cite[6.19]{Taylor2014Maximal}. That image has square one
\cite[Proposition~4.11(a)]{Bonnafe2005}.
Since $Z(\mathbf H^*)$ is a torus, choose
$z\in Z(\mathbf H^*)$ with $(zs_0)^2=1$ and set $\widetilde s=zs_0$.
Multiplication by $z$ leaves the centraliser and its Weyl group family
unchanged. Since $\widetilde s^2=1$ and $q_0$ is odd, $F_0$ fixes
$\widetilde s$ in this split torus and acts trivially on $W_{\widetilde s}$.
The ordinary family used below is
$\Irr_{\widetilde s,\mathcal F}(H)$ for this rational pair.

The construction in \cite[Proposition~10.2]{Taylor2014Maximal}
identifies truncated induction of the special character of
$\mathcal F$ with the Springer character of the chosen unipotent
class $\mathcal C$. The family group and the component group
$A_{\mathbf H}(u)=C_{\mathbf H}(u)/C_{\mathbf H}(u)^\circ$, for
$u\in\mathcal C$, are elementary abelian of the same order, as
established in the manuscript using
\cite[7.1 and Proposition~10.2]{Taylor2014Maximal}.
The argument in \cite[7.4]{Taylor2014Maximal} applies to
$\widetilde s^2=1$ and gives condition~\textup{(P4)} of
\cite[Theorem~2.11]{Taylor2014Maximal}. This supplies the geometric
hypothesis \cite[(4.4)(b)]{Geck1999GGGR} used in place of isolation.

The proof, in Appendix B of the manuscript, applies the geometric argument of Lemma 2.11
to these Type C families. An auxiliary Frobenius $F'=F_0^{2N}$,
with $N$ sufficiently divisible, satisfies
\cite[(4.7)(a)--(c)]{Geck1999GGGR} and has field size
$q'=q_0^{2N}\equiv1\pmod4$. The diagonal and support identities,
the comparison of Fourier constants and the multiplicity calculation
establish \cite[Proposition~3.1 and Corollary~3.2]{Geck1999GGGR}
for the groups required in the argument. The constants needed for
the restriction theorem then follow from
\cite[Theorem~3.8]{Geck1999GGGR}.
The resulting restriction multiplicities are independent of the
rational structure. For the original $F_0$, the local system trace
values give the character table of $A_{\mathbf H}(u)$ up to nonzero
row factors. Shoji's theorem \cite[Theorem~3.2]{Shoji1995II} expresses
the corresponding character sheaf trace functions as nonzero scalar
multiples of almost characters in the same ordinary family. Thus
ordinary values from that family span all functions on the rational
classes in $\mathcal C^{F_0}$.

For this original Frobenius, triviality of its action on
$W_{\widetilde s}$ and the abelian family group give the generic
denominator formula \cite[(6.1)(c)]{Geck1999GGGR}.
Together with the weighted scalar product formula and nonnegative
integral multiplicities, it makes every scalar product row have
sum one. The spanning and independence arguments show that every
column has a nonzero entry. This gives the required Kronecker delta
scalar products with the generalised Gelfand--Graev characters.

Proposition 3.3 descends this selection to $G_0=\Sp_{2m}(q_0)$.
The nondegenerate symbols in each type $\mathsf D$ factor of the
selected family allow the argument of \cite[10.3]{Taylor2014Maximal}
to give condition~\textup{(P5)} of \cite[Theorem~3.2]{Taylor2013}
for the same split pair. The restriction argument in
\cite[Proposition~5.4]{Taylor2013} uses this condition, induction of
generalised Gelfand--Graev characters
\cite[proof of Lemma~14.12]{Taylor2016}, compatibility with
Alvis--Curtis duality and Frobenius reciprocity. It relates the
restriction constituents to the rational classes of $G_0$ in each
$H$-class. Their unipotent supports are preserved under restriction
\cite[Example~2.7.18]{GeckMalle2020}. Their rational Lusztig series
has parameter the image of $\widetilde s$
\cite[Proposition~2.6.16]{GeckMalle2020}. This image has square one,
so the constituents and their positive Alvis--Curtis duals belong
to the principal block
\cite[Proposition~9.8(iv) and Theorem~21.14]{CabanesEnguehard2004}.
The vanishing results for wave front sets give the zero entries needed for unitriangularity in \cite[Theorem~2.9]{Chaneb2021}. This is the
character selection used in the principal block argument of
\cite[Section~4.1 and Corollary~4.3]{FengMalle2022}.

The total numbers and the numbers fixed by the diagonal action follow by
combining the Brauer character parametrisation with the counting identities in
\cite[Corollary~4.3, Proposition~6.1 and proof of Theorem~6.2]
{FengMalle2022}. Inner automorphisms and field automorphisms act trivially
on the sets being compared. The remaining diagonal action is involutory on
those sets. The general fixed point count argument in the library manual
therefore constructs the equivariant bijection.

\begin{implementation}
The main files are
\module{ManuscriptIBAW.TypeC.PrincipalConverse},
\module{ManuscriptIBAW.TypeC.PrincipalProducts},
\module{ManuscriptIBAW.TypeC.PrincipalParameters},
\module{ManuscriptIBAW.TypeC.PrincipalNumericalParameters},
\module{ManuscriptIBAW.TypeC.PrincipalParameterCoordinates} and
\module{ManuscriptIBAW.TypeC.PrincipalParameterActions}.
The structures \lean{PrincipalFactorData} and
\lean{PrincipalFactorSelection} in \lean{PrincipalParameters} specify
the factor decompositions used by the inverse construction.
\module{ManuscriptIBAW.TypeC.PrincipalEquivariance} constructs the
correspondence from the specified actions and counts.
In \module{ManuscriptIBAW.TypeC.GGGRSelection},
\lean{delta_selection} constructs distinct characters in the specified
ordinary family of $\CSp_{2m}$ whose positive Alvis--Curtis duals have
the required Kronecker delta scalar products. In the application,
this is the family indexed by the centrally normalised split pair
$(\widetilde s,\mathcal F)$ above. The construction of that pair,
condition~\textup{(P4)} and the geometric interpretation of the
family are not formalised. The scalar products are normalised by
the group order and use inverse character values, with trace
identities in the same convention.
The trace expansions for the original Frobenius, support and
duality identities, nonnegative multiplicities, weighted formula,
denominator equality and trivial component action are explicit
assumptions. Their geometric proofs, including the auxiliary
Frobenius argument, remain outside the formalisation.
Lean derives spanning, independence, localisation of scalar
products and a nonzero entry in every column. It also proves
that the component indices are one and every row has sum one.

\module{ManuscriptIBAW.TypeC.PrincipalGGGR} selects characters of
the symplectic group using the specified restriction identities,
Frobenius reciprocity and Clifford action. In \lean{P5Source},
\lean{p5_count} assumes the following numerical consequence of
condition~\textup{(P5)} and the class count in
\cite[Proposition~5.4]{Taylor2013}: whenever the scalar product on $H$ is nonzero, the number of constituents on restriction to $G_0$ equals the number of $G_0$-conjugacy classes in the corresponding $H$-conjugacy class.
Lean proves the column sums from the restriction identities.
Clifford transitivity and the conjugation identities make all row
sums equal, and the assumed count forces them to be one.
This selects the required characters of $G_0$. Geometric support and
wave front vanishing remain explicit assumptions.
The finite character selection is a separate result. Its interpretation
as ordinary characters in the principal block, and the deduction of
principal block Brauer character counts from the selected rows, are
not formalised. The principal block application instead assumes
\lean{FengMalleCorollary43aSource} and \lean{FMPrincipalCountSource}
directly, as the fields \lean{fieldFixed} and \lean{counts} of
\lean{PrincipalGGGR.PrincipalSource}. These are the field invariance
and counting conclusions used from Feng--Malle. The application does
not derive them from \lean{lowerTriangular}.

The application to all blocks separately assumes
\cite[Corollary~4.6]{FengMalle2022} as
\lean{LiteralFengMalleCorollary46Certificate}.
It is not a consequence of the principal character selection.
The covering and character triple hypotheses also remain explicit.

For the application to all blocks in rank two,
\module{ManuscriptIBAW.TypeC.LiLiSource} takes Li--Li's theorem as an
external result for each block at a nondefining prime over an odd field,
in the specified splitting modular system. Its application requires
an identification of the acting group with the full automorphism
stabiliser of that block. This identification is the hypothesis
\lean{Definition35AutomorphismStabilizerAdapter}, and is not a
conclusion asserted for an arbitrary group action. The cover is the identity on $\PSp_4$ at two and the map
$\Sp_4\to\PSp_4$ at odd primes.
\module{ManuscriptIBAW.TypeC.LowRankApplications} combines these
blockwise conclusions using a separate compatibility assumption across
blocks.

The Type C principal family uses an algebraically closed ordinary
coefficient field. The source \lean{ClosedSource} interprets Li--Li's
theorem directly in this field, with the block and root conventions in
\lean{ClosedCoefficients}. This interpretation is an external assumption.
\lean{ClosedRankTwoPrincipalInputs.seed} lifts the resulting quotient
block bijection to $\Sp_4$. For each Brauer character and its corresponding
weight, a separate assumption identifies the central
quotients defined by the reference character and the given Brauer character
and supplies compatible roots of unity, extensions and intermediate block
data. The assumptions \lean{SingleBlockCompatibilitySource} and
\lean{ClosedSingleBlockCompatibilitySource} apply to every coefficient
prime, family, cover and coefficient choice of the specified types,
every block and block witness, and every Brauer character in that block.
They are not restricted to rank two and concern the specified universal
$\ell'$-covers. Each supplied block witness already identifies its acting
group with the full automorphism stabiliser of the block.
These applications use the instances selected by the central quotient
construction. The given bijection and its equivariance are preserved.
The central lift retains its inflation and character triple hypotheses.
\end{implementation}

For the modular character triple and covering arguments, choose the local
characters in a common root convention. An adjustment within diagonal orbits
gives a correspondence satisfying the required triple relation for every
character and its weight. The extension and covering steps use
\cite[Proposition~3.4, Remark~3.5 and Corollary~4.6]{FengMalle2022} with
those compatible choices.

\section{The unipotent correspondence over even fields}

For even $q$, Li \cite[Section~6]{Li2021} obtains the descriptions of
blocks and weights by adapting arguments of An and Fong--Srinivasan.
Li proves the inductive condition under a unitriangularity assumption on the relevant decomposition matrix \cite[Theorem~3]{Li2021}. In rank at least four, the manuscript
gives an independent proof using generic weights and Jordan reduction,
without that assumption.

Let $H=\Sp_{2r}(2^a)$ with $a>0$, let $\ell$ be odd, and let $e$ be the
multiplicative order of $2^a$ modulo $\ell$. For a unipotent block $b$, let
$X_b$ be its unipotent ordinary characters and let $\mathcal W(b)$ be the
set of conjugacy classes of generic weights belonging to $b$.

The unipotent block parametrisation and the relative Weyl group argument
must use the same pair. Cabanes--Enguehard's Theorem 4.4 identifies
the relevant series \cite{CabanesEnguehard1994}. We use the relative Weyl
group parametrisation recalled in Cabanes--Sp\"ath's Theorem 2.8
\cite{CabanesSpath2013}. Kessar--Malle's Definition 2.1, Remark 2.2(2) and
Theorem A give the comparison with the required $e$-Jordan-cuspidal pair in
this setting \cite{KessarMalle2015}. The uniqueness and block membership
statements identify this pair and establish that it is
$e$-Jordan-generalised-cuspidal.

Theorem 7.5 of Feng--Malle--Zhang then gives the relative Weyl group
correspondence \cite{FengMalleZhang2026}. Their Lemma 3.21 concerns one
pair. Passing to all generic weights also uses Definitions 3.17--3.18, the
conjugacy classes of pairs and the character correspondence on their
inertia quotients. The passage to all weights is conditional on the
required extendibility for every pair, and the admissible pairs are
explicitly $e$-split. The resulting character correspondence also gives the
cardinality equality.

The field invariance argument uses a rational Levi parametrisation and the
same Lang conjugator throughout. Let $\mathbf M$ be the selected rational
Levi subgroup and let $F'$ define $H$. The field automorphism of a standard
Levi is transported to an automorphism $\tau$ fixing the selected pair. The
action of $\tau$ on $N_H(\mathbf M)/\mathbf M^{F'}$ is trivial. Extend the
unipotent character to the required stabiliser using
\cite[Theorem~3.3]{Spath2025Extensions}. This application uses the embedding of $\mathbf H$ into the simple adjoint group $\mathrm{PGL}_{2r+1}$ that sends $A$ to the projective image of $\operatorname{diag}(A,1)$. The extension group is
the simultaneous stabiliser of the subgroup and its unipotent character,
including the prescribed field automorphisms. The geometric hypotheses
identify this stabiliser and the inclusion of the finite Levi subgroup.
Gallagher's theorem and induction then give invariance of the generic weight
class. The ambient Frobenius powers are required to induce the specified
automorphisms on the finite fixed point groups. No corresponding finite
order relation on the entire algebraic group is assumed.

The integral basic set in \cite[Theorem~A]{Geck1993} identifies the free
abelian groups on $X_b$ and $\IBr(b)$ through the specified decomposition
homomorphism. Field automorphisms fix $X_b$ pointwise, so the elementary
basis argument proves that they also fix $\IBr(b)$ pointwise. Equality of
the ranks gives an equivariant bijection between these sets.

The complex ordinary characters and the ordinary labels in the basic set are
related through a common cyclotomic field. The two fixed embeddings identify
the values of the same character and therefore determine the automorphism
action on its label over the ordinary
coefficient field. The comparison of labels thus commutes with the given
automorphisms.

Under Condition 6.1, Theorem 6.2 of \cite{FengMalleZhang2026} gives an
equivariant correspondence between generic weights and Alperin weights,
preserving blocks. Here $\ell$ is good and does not divide
$2\cdot2^a|Z(H)|$. The regular embedding and automorphism group are those
specified in that condition. Manuscript Lemma 2.10 applies the specialised
Brough--Sp\"ath criterion, in the form stated in
\cite[Theorem~3.18]{FengLiZhang2022Jordan}, to show that the block is
BAW-good. The passage from BAW-goodness to the modular character triple
condition uses the stated identifications of the local character reductions
and compatible roots.

\begin{implementation}
\module{ManuscriptIBAW.TypeC.PairLabelComparison} identifies the initial
pair. \module{ManuscriptIBAW.Characters.CyclotomicOrdinaryLabels}
compares character values over the two coefficient fields. The direct basic
set and unipotent arguments are in \module{ManuscriptIBAW.TypeC.EvenBasicSet} and
\module{ManuscriptIBAW.TypeC.EvenUnipotent}.
The local semisimple label comparison in manuscript Lemma 3.6 is an explicit external
assumption, as are the rational Levi parametrisation and the extension,
Clifford and Gallagher formulas. Given these statements, the formal proof
derives field invariance of the generic weight classes. It does not formalise
the Deligne--Lusztig theory underlying the label comparison.
The modules beginning
\lean{ManuscriptIBAW.TypeC.EvenApplication} apply the criterion to the specified groups and covering maps under the stated source hypotheses.
\end{implementation}

\section{Jordan reduction for each semisimple parameter}

Jordan reduction uses every pair $(\mathbf K,F')$ required by Feng--Li--Zhang's Hypothesis 5.5 \cite{FengLiZhang2022Jordan}, including the pair to which the reduction is being applied. The corresponding group $\mathbf K^{F'}$ is included even if it has a nontrivial central $\ell$-subgroup or is not itself a universal $\ell'$-cover.

The automorphism group in the stabiliser and extension hypothesis is chosen
separately for each semisimple parameter $s$, in accordance with the
construction in Section 5 and Remark 5.4 of the same paper. Restrict the
stabiliser factorisation and a character extension to the chosen subgroup,
retaining the representation affording the same character. Irreducibility
after restriction follows because the restricted group still contains the
base group on which that representation is irreducible.

The geometric choices identify field automorphisms representing the required
outer automorphisms stabilising the series of $s$. Compatible Frobenius
representatives and a Levi subgroup in duality with the minimal dual Levi
are chosen by the construction using inner conjugacy in
\cite[Corollary~4.2, Lemma~4.5 and Proposition~4.6]{Ruhstorfer2022Derived},
preserving the same series idempotent and the required outer stabiliser.
These choices do not assert stability of an arbitrary preselected Levi.
The stabiliser and extension hypotheses for this chosen group are
established before applying Theorem 5.7.

The geometric series use the original blocks. If $t_b$ is the
semisimple $\ell'$-parameter assigned to $b$, the series idempotent at $s$
is identified with the sum of the primitive idempotents whose parameters are
conjugate to $s$ in the finite dual group. Orthogonality proves that a block
belongs to the series exactly when this conjugacy condition holds. The
specified inclusion of finite rational points in the algebraic dual
sends this same $t_b$ to the parameter used in the strict quasi-isolation
argument. Thus the block of the lower group and its geometric series refer
to the same parameter before the reduction theorem is applied. The
interpretation of the supplied series, dual group and semisimple elements
remains an external hypothesis.

For odd $q$ and $\ell=2$, the applicable quasi-isolation statement forces the
semisimple odd-order label to be the identity, and the lower symplectic
block is principal \cite[Proposition~4.11(a)]{Bonnafe2005}.
For characteristic two and odd $\ell$, the quasi-isolated label is also
the identity, so the block is unipotent \cite[Example~4.8]{Bonnafe2005}. These applications
require the stated semisimplicity, dual group type and characteristic
conditions.

For a rank two relative factor over an odd field, the manuscript uses
the principal $\Sp_4$ block bijection obtained from Li--Li's theorem by
the central lift in Proposition 3.3. For higher rank factors
it uses the principal construction above.

The lower type A cases use the special linear or unitary fixed point
groups, including their central $\ell$-subgroups when present. The invoked
result retains the type A hypotheses
\cite{FengLiZhang2023Morita} and those for the exceptional small
groups \cite{FengLiZhang2023LowRankA}. For split low ranks
over even fields of order at least four, Schaeffer Fry's Theorem 5.5
applies. For $\Sp_6(2)$ it applies at every odd prime on the specified
double cover \cite{SchaefferFry2014}. The Suzuki cases use
\cite[Corollary~6.3]{Spath2013}, together with Proposition~4.6 of the same
paper for the required descent to the fixed point group. The specified
covers are retained, including the exceptional cover for ${}^2B_2(8)$.

\begin{implementation}
The general restriction and geometric constructions are in
\lean{ManuscriptIBAW.Jordan}, including the block and label
identities in \module{ManuscriptIBAW.Jordan.SeriesBinding}.
The specialised applications are
\module{ManuscriptIBAW.Jordan.TypeCOddLocalized} and
\module{ManuscriptIBAW.Jordan.TypeCEvenLocalized}.
\module{ManuscriptIBAW.TypeC.PrincipalApplicationComplete} and
\module{ManuscriptIBAW.TypeC.LowRankApplications} keep the corresponding
complete family and the specified groups in small rank.
Rank two principal relative factors use the specified Li--Li application
and central lift. Higher rank factors use the principal construction
with its stated character theoretic assumptions.
\end{implementation}

\section{Odd nondefining primes over odd fields}

Let $\widetilde G=\CSp_{2n}(F)$, let $E$ be the cyclic group of field
automorphisms, and let $C$ be the group of ordinary linear characters of
$\widetilde G/G$ of order prime to $\ell$. The multiplier identifies
$\widetilde G/G$ with $F^\times$ and sends the centre onto its squares. Thus
$|\widetilde G:GZ(\widetilde G)|=2$ when $q$ is odd.

For a block stabiliser $J$ in $C\rtimes E$, the kernel $D$ of the field
projection embeds in its linear character stabiliser. The general central
character argument gives $|D|\leq2$, and every subgroup of $J$ is
$2$-hypoelementary. The multiplier and central character calculations
therefore give both the quotient index and the block stabiliser bound.

The required integral basic set is given by
\cite[Theorem~5.1]{GeckHiss1991}. Its decomposition homomorphism, root
convention and actions remain specified. Conlon's theorem \cite{Conlon1968},
in the formulation of \cite[Theorem~3.5.5]{Bouc2000}, gives the equivariant
comparison needed here. The prime used for the permutation lattice is
$2$, while $\ell$ is odd.

For the symplectic group, the ordinary basic set is given by Theorem A
\cite{Geck1993}. Here $\mathcal E_{\ell'}(G)$ is defined from the rational
Lusztig partition on complex characters by requiring a semisimple dual
parameter of order prime to $\ell$. The dual groups are the split special
orthogonal group for $G$ and the special Clifford group for $\widetilde G$.
Fixed embeddings of one cyclotomic field compare the same character values
over $\C$ and the ordinary coefficient field. The sets used in the two
basic set arguments are assumed to be these unions of rational series. Their identification with the Lusztig partition in the
cited results remains an explicit hypothesis.

Cabanes--Enguehard Proposition 15.6(i) gives the restriction comparison for
each fixed rational semisimple lift of a parameter from the special orthogonal dual to the special Clifford dual
\cite{CabanesEnguehard2004}. Membership in the lower series is equivalent to
occurrence in the restriction of an upper character in that fixed upper
series. A separate structural assumption supplies the semisimple lift under
the specified projection from the special Clifford group to the special
orthogonal group. Reindexing the ordinary restriction scalar product proves
that conjugation by $\widetilde G$ preserves occurrence as a restriction constituent. The same upper lift and upper
character then prove diagonal stability of the lower series. The lower
parameter is unchanged, so its order remains prime to $\ell$. This argument
does not require an upper lift of order prime to $\ell$. The source
interpretation uses the specified regular embedding, equality of derived
groups, connected centre of the conformal algebraic group and connected central dual kernel.

For the field action, let $\sigma$ be the specified field automorphism of $G$,
and let $\sigma^*$ denote the specified coordinate field automorphism of the
split orthogonal dual. Taylor's Proposition 7.2 transports $\chi$ by
pullback $\chi\mapsto\chi\circ\sigma^{-1}$ and the label $s$ by $(\sigma^*)^{-1}(s)$
\cite{Taylor2018Automorphisms}. The algebraic dual isogeny identification
and the common torus and root conventions remain part of this ordinary
source statement. The inverse dual automorphism preserves order, so
semisimplicity and order prime to $\ell$ are preserved. The common
cyclotomic character identities give the same conclusion over the ordinary
coefficient field.

These calculations prove stability of the union under the full conformal
and field action. Field automorphisms may permute its individual rational
series. On this union, the remaining basic set assumptions give the integral
decomposition homomorphism and its compatibility with blocks. Apply the
block basic set argument over representatives of block orbits and transport
the resulting bijections over their orbits.
This constructs one bijection preserving blocks
\[
 \beta:\IBr(G)\longrightarrow\mathcal E_{\ell'}(G)
\]
commuting with the conformal and field actions. Surjectivity of their
action onto $\Aut(G)$ gives equivariance under every automorphism.

The ordinary character and weight formulae in
\cite[proofs of Lemmas~5.1 and~5.3]{Li2021} describe the same changes of
labels under tensoring and field automorphisms. The unitriangularity
assumption in those lemmas is used to transfer ordinary labels to Brauer
characters. The comparison of ordinary character and weight labels used
here does not require that assumption.
The ordinary stabiliser factorisation is supplied by
\cite[Theorem~3.1]{CabanesSpath2017TypeC}.

Apply that ordinary theorem to $\beta(\varphi)$. Equivariance and
injectivity of $\beta$ transfer the factorisation of conformal and field
stabilisers to every $\varphi\in\IBr(G)$. The common character identities
transport the ordinary theorem from $\C$ to the coefficient field. The
restriction expansion can then select a constituent with the required
stabiliser.

Li proves the weight stabiliser factorisation for the standard radical
subgroups described in \cite[Section~2.C]{Li2021}, in
\cite[Lemma~5.5(1)]{Li2021}. Conjugating the chosen representative gives
the factorisation for every weight. The manuscript also uses
\cite[Lemma~5.5(2)]{Li2021} for the ordinary character extension in
condition~(iv) of \cite[Theorem~4.5]{BroughSpath2022}.

The character and weight stabiliser factorisations, cyclic extensions and
subgroup index calculations give the hypotheses of the criterion in
\cite[Theorem~4.5]{BroughSpath2022} on the specified symplectic cover.
The passage to the complete family condition also uses
\cite[Lemma~3.3]{KoshitaniSpath2016} and the normalisation in
\cite[Definition~4.1]{Spath2013}, with the specified roots and block
operations.

\begin{implementation}
The conformal correspondence is
\module{ManuscriptIBAW.TypeC.OddConformalCorrespondence}.
\module{ManuscriptIBAW.TypeC.OddRationalSeries} defines the two rational
unions. \module{ManuscriptIBAW.TypeC.OddRationalAutomorphisms}
proves stability of the symplectic union under diagonal and field
automorphisms from the stated ordinary character formulae.
\module{ManuscriptIBAW.TypeC.OddLiteralBasicSet}
combines this stability with the integral decomposition hypotheses.
The assumptions are given by
\lean{OddPrimeApplication.SourceInputs.basicSetData}, from which
\lean{basicSet} is constructed.
\module{ManuscriptIBAW.TypeC.OddGlobalFactorization} constructs
the global bijection and transfers the stabilisers.
\module{ManuscriptIBAW.TypeC.OddPrimeApplication} derives the index
consequences and applies the criterion. Its field
\lean{SourceInputs.li} has type
\lean{TypeCCurrentOddPrimeLiSource.Li55SplittingCertificate}. This
external assumption requires the standard radical subgroup family and
its weight stabiliser factorisation over every coefficient field of
characteristic zero containing the required roots of unity for the full
semidirect product.
The ordinary character extensions use separate cyclic extension
assumptions on the two inertia quotients. The Brauer character extension
uses a separate assumption.
\module{ManuscriptIBAW.TypeC.HigherOddApplication} specifies the
cover and simple quotient before invoking that result.
\end{implementation}
